\subsection{Introduction \quad / \quad \textthai{บทนำ}}
\begin{paracol}{2}
In the history of classical mechanics, two frameworks stand out as the most powerful and elegant: the Lagrangian and the Hamiltonian. These frameworks move beyond Newton's "Force" view to an "Energy" view of the universe. By focusing on scalars like kinetic and potential energy, we can describe complex systems with relative ease. However, when we apply standard Artificial Intelligence (like basic MLPs) to these systems, we often find that the models "violate" physics—they fail to conserve energy and suffer from numerical drift.

This chapter introduces \textbf{Hamiltonian Neural Networks}\index{HNNs@HNNs (HNNs)} (HNNs), a class of physics-inspired architectures that respect the underlying \textbf{Symplectic Geometry}\index{Symplectic Geometry@เรขาคณิตแบบซิมเพล็กติก (Symplectic Geometry)} of classical mechanics. By learning the energy function directly, HNNs ensure that simulations are physically consistent even over long time horizons.

\switchcolumn

\begin{thai}
ปัญญาประดิษฐ์แบบมาตรฐานมักมองข้าม \textbf{กฎการอนุรักษ์ (Conservation Laws)}\index{Conservation Laws@กฎการอนุรักษ์ (Conservation Laws)} ทำให้แบบจำลองทางฟิสิกส์ผิดเพี้ยนไปเมื่อเวลาผ่านไป Hamiltonian Neural Networks (HNNs) แก้ปัญหานี้โดยการเรียนรู้ "ฮามิลโทเนียน" ของระบบแทนการเรียนรู้ความเร่งโดยตรง ทำให้ AI เคารพกฎการอนุรักษ์พลังงานและความสมมาตรของฟิสิกส์ ผ่านพื้นฐานของ\textbf{กลศาสตร์ลากรานจ์ (Lagrangian Mechanics)}\index{Lagrangian Mechanics@กลศาสตร์ลากรานจ์ (Lagrangian Mechanics)} และฮามิลโทเนียน
\end{thai}
\end{paracol}
